%
% hoehe.tex -- template for standalon tikz images
%
% (c) 2021 Prof Dr Andreas Müller, OST Ostschweizer Fachhochschule
%
\documentclass[tikz]{standalone}
\usepackage{amsmath}
\usepackage{times}
\usepackage{txfonts}
\usepackage{pgfplots}
\usepackage{csvsimple}
\usetikzlibrary{arrows,intersections,math,calc}
\begin{document}
\def\skala{1}
\begin{tikzpicture}[>=latex,thick,scale=\skala]

\coordinate (A) at (0,0);
\coordinate (B) at (5,0);
\pgfmathparse{sqrt((5-1.8)*1.8)}
\xdef\h{\pgfmathresult}
\coordinate (C) at (1.8, \h);
\coordinate (H) at (1.8, 0);
\node at (A) [left] {$A$};
\node at (B) [right] {$B$};
\node at (C) [above] {$C$};
\draw (A) -- (B) -- (C) -- cycle;
\draw[dashed] (C) -- (H);
% Rechter Winkel bei H
\draw (H) ++(0.18,0) -- ++(0,0.18) -- ++(-0.18,0);
% Seiten-Labels klar ausserhalb
\node[left=6pt]        at ($0.5*(A)+0.5*(C)$) [above left] {$b$};
\node[right=6pt]       at ($0.5*(B)+0.5*(C)$) [above right] {$a$};
\node[below=8pt]       at ($0.5*(A)+0.5*(B)$) [below] {$c$};
\node[right=6pt]       at ($0.5*(H)+0.5*(C)$) [right] {$h$};
% Gleichung unterhalb des Dreiecks, klar getrennt
\node[below=18pt] at (2.5, 0) {$\dfrac{1}{2}ch \;=\; \dfrac{1}{2}ab$};

\end{tikzpicture}
\end{document}

